\section{Evaluation}

The evaluation of ECHOREPO focuses on its behaviour as an operational
research-data infrastructure rather than on isolated software benchmarks.
Four aspects are considered: the scale of the production deployment,
the ability of the quality-assurance workflow to identify and recover
problematic citizen-generated records, the generation and external
publication of reusable research products, and the computational
requirements associated with the different data domains.

\subsection{Operational deployment scale}

ECHOREPO has been operated as the research-data infrastructure supporting
the ECHO soil citizen-science campaign. At the time of writing, the
production repository contains around 9,000 soil samples collected
across multiple European countries and manages approximately 100~GB of
data.

The different data domains have markedly different storage and processing
requirements. Sample-level observations and laboratory parameters are
relatively compact structured records, while images contribute larger
binary objects. Microbial biodiversity is the dominant component in terms
of both storage volume and computational requirements because individual
samples may be associated with large numbers of taxonomic features and
abundance measurements.

Table~\ref{tab:deployment-scale} summarises the principal characteristics
of the deployment.

\begin{table}[htbp]
    \centering
    \caption{Scale of the operational ECHOREPO deployment at the time of
    evaluation. Values marked TODO will be generated directly from the
    production repository before submission.}
    \label{tab:deployment-scale}
    \begin{tabular}{lr}
        \toprule
        Metric & Value \\
        \midrule
        Soil samples & $>10{,}000$ \\
        Countries represented & [TODO] \\
        Sample images & [TODO] \\
        Laboratory parameter records & [TODO] \\
        Samples with biodiversity data & [TODO] \\
        Biodiversity features & [TODO] \\
        Non-zero biodiversity abundance records & [TODO] \\
        Total managed data volume & $\sim100$~GB \\
        \bottomrule
    \end{tabular}
\end{table}

These figures are relevant to the architectural evaluation because the
repository must accommodate substantially different workloads within the
same sample-centred data lifecycle. In particular, the biodiversity
component demonstrates that enrichment of an initially lightweight
citizen-science observation can eventually produce datasets whose
storage and processing requirements are substantially larger than
the original field record.

\subsection{Quality assurance and data recovery}

The geographical quality-control analysis described in
Section~\ref{sec:data-quality} provides the clearest quantitative
evaluation of the QA workflow.

In the analysed snapshot of 6,871 samples, 492 records (7.16\%) were
classified as having problematic coordinates. These included 281 cases
in which a country could not be determined from the recorded position,
123 cases in which the inferred country did not correspond to the
expected country, and 88 records containing known default coordinates.
A further three records lacked the reference country information required
for the corresponding validation procedure.

The significance of these results is not limited to the error rate itself.
They demonstrate that, even when a data-collection application performs
validation at submission time, additional repository-level quality
control remains necessary once observations are considered in their
broader geographical and project context.

The second component of the evaluation concerns the ability to recover
records after an error has been detected. ECHOREPO retains sample
identifiers, participant associations and other provenance information
that can be used to reconstruct or correct information in cases where
sufficient evidence is available.

The final evaluation will distinguish the outcomes shown in
Table~\ref{tab:recovery-results}.

\begin{table}[htbp]
    \centering
    \caption{Planned evaluation of recovery outcomes for records flagged
    by the quality-assurance workflow.}
    \label{tab:recovery-results}
    \begin{tabular}{lrr}
        \toprule
        Outcome & Records & Percentage of flagged records \\
        \midrule
        Automatically recovered & [TODO] & [TODO] \\
        Recovered after review & [TODO] & [TODO] \\
        Retained but still flagged & [TODO] & [TODO] \\
        Unresolved / unusable & [TODO] & [TODO] \\
        \bottomrule
    \end{tabular}
\end{table}

This distinction is important because the performance of the QA system
cannot be characterised solely by the number of errors it detects. From
a research-data perspective, the more relevant question is how many
scientifically useful records can ultimately be retained after the
problem has been identified.

The recovery rate will therefore be calculated as:

\begin{equation}
    R =
    \frac{N_{\mathrm{recovered}}}
         {N_{\mathrm{flagged}}}
    \times 100\%,
\end{equation}

where $N_{\mathrm{flagged}}$ is the number of records identified as
problematic and $N_{\mathrm{recovered}}$ is the subset for which a
sufficiently supported correction can be established.

\subsection{Publication and external interoperability}

A further evaluation criterion is whether ECHOREPO can transform its
operational data into research products that remain usable independently
of the running application.

The canonical publication workflow currently generates separate
machine-readable resources describing samples, images, analytical
parameters and biodiversity results, together with supporting definitions
of parameters and analytical methods. The principal canonical resources
are:

\begin{itemize}
    \item \texttt{samples.csv};
    \item \texttt{sample\_images.csv};
    \item \texttt{sample\_parameters.csv};
    \item \texttt{sample\_biodiversity.csv};
    \item \texttt{parameter\_definitions.csv}; and
    \item \texttt{analysis\_methods.csv}.
\end{itemize}

These resources can be generated from the operational repository and
retrieved through machine-oriented interfaces without relying on the
interactive ECHOREPO web application.

The publication workflow has additionally been exercised through Zenodo.
Published releases receive persistent identifiers and form a versioned
publication history, allowing a particular research-data release to be
cited independently of the continuously evolving production database.
The deposited resources are accompanied by structured metadata intended
to support machine interpretation and external harvesting.

The resulting interoperability chain has been demonstrated across
multiple independently operated systems:

\begin{verbatim}
ECHOREPO
    |
    v
canonical research resources
    |
    v
Zenodo publication
    |
    v
persistent identifier
    |
    v
SoilWise discovery
\end{verbatim}

The ability of a dataset originating in ECHOREPO to be deposited in an
external trusted repository and subsequently discovered through SoilWise
provides an end-to-end test of the publication workflow. It demonstrates
that access to the research products is not dependent on knowledge of
ECHOREPO's internal database or web interface.

For the final manuscript, this evaluation will record the corresponding
Zenodo concept DOI, the publication version used for the evaluation, and
the associated SoilWise catalogue record.

\subsection{Operational performance and scalability}

For the final manuscript, we also plan to characterise representative
operational performance and scalability.

ECHOREPO has been developed primarily to support an operational research
workflow rather than as a high-performance data-processing platform.
Nevertheless, representative processing measurements are useful for
understanding the infrastructure required to reproduce the deployment.

The final evaluation will therefore report representative execution times
for several complete operations rather than micro-benchmarks of individual
software components. These measurements will include:

\begin{itemize}
    \item ingestion and canonicalisation of citizen-science records;
    \item generation of the complete canonical publication bundle;
    \item ingestion or processing of a representative biodiversity dataset;
    \item generation of biodiversity publication resources; and
    \item the approximate storage contribution of the main data domains.
\end{itemize}

Measurements will be performed on the production or a representative
deployment and will report the relevant hardware and software environment.
Their purpose is not to compare ECHOREPO with alternative database
technologies, but to provide prospective adopters with an indication of
the resources required for deployments of comparable scale.

\begin{table}[htbp]
    \centering
    \caption{Representative operational performance measurements TO BE COMPLETED BEFORE SUBMISSION.}
    \label{tab:performance}
    \begin{tabular}{lrr}
        \toprule
        Operation & Data volume & Execution time \\
        \midrule
        Canonical sample export & [TODO] & [TODO] \\
        Complete canonical bundle & [TODO] & [TODO] \\
        Biodiversity ingestion & [TODO] & [TODO] \\
        Biodiversity export & [TODO] & [TODO] \\
        \bottomrule
    \end{tabular}
\end{table}

Taken together, these evaluations address different aspects of the
infrastructure. Deployment scale establishes that ECHOREPO has been used
beyond a prototype setting; the QA analysis evaluates its handling of
imperfect citizen-generated data; the publication workflow demonstrates
that internal records can be transformed into independently discoverable
research products; and the operational measurements provide information
about the resources required to reproduce the approach.